\documentclass[]{article}

\usepackage{amsmath}
\usepackage{multirow}
\usepackage{natbib}
\usepackage{graphicx} 


\begin{document}

Galaxy clusters, as the most massive gravitationally bound structures in the cosmos, provide a crucial laboratory for studying cosmology and the formation of large-scale structure. The abundance and spatial distribution of these objects are highly sensitive to fundamental cosmological parameters, such as the total matter density and the amplitude of primordial density fluctuations. The thermal Sunyaev-Zel'dovich (tSZ) effect, a spectral distortion of the Cosmic Microwave Background (CMB) caused by inverse Compton scattering of photons off hot electrons in the intracluster medium, offers a robust method for their detection. A key advantage of the tSZ effect is that its surface brightness is nearly independent of redshift, enabling the construction of large, mass-selected cluster catalogs extending to the early universe. Upcoming surveys are expected to detect tens of thousands of clusters using this technique, promising to significantly refine our cosmological model.

The primary challenge in realizing the full scientific potential of these surveys lies in cleanly extracting the faint tSZ signal from millimeter-wave sky maps. These maps contain a superposition of signals, including the dominant primary CMB anisotropies and various astrophysical foregrounds. While the primary CMB can be effectively separated due to its distinct blackbody spectrum, a more pernicious contaminant is the Cosmic Infrared Background (CIB), which arises from the integrated emission of dusty, star-forming galaxies. The CIB is a significant source of emission at the very frequencies where tSZ experiments are most sensitive. Critically, because these galaxies inhabit the same dark matter halos as galaxy clusters, the CIB is spatially correlated with the tSZ signal itself. This correlation complicates separation based on angular morphology alone and can introduce significant biases in the reconstructed tSZ map, leading to spurious cluster detections and an increase in the scatter of the mass-observable relation.

A widely used technique for this component separation task is the Internal Linear Combination (ILC) method. The ILC constructs a weighted sum of multi-frequency maps, with weights chosen to minimize the total variance in the output map while preserving signals with a specific spectral signature, such as that of the tSZ effect. However, standard ILC implementations can be suboptimal. They often treat foregrounds as contaminants to be deterministically nulled, a constraint that can lead to an undesirable amplification of instrumental noise. This occurs when frequency channels with high foreground contamination but low intrinsic noise are down-weighted, forcing the algorithm to rely more heavily on noisier channels to enforce the null condition. This creates a difficult trade-off between foreground removal and noise amplification.

In this paper, we implement a more statistically robust framework based on a Multi-Frequency Wiener Filter (MWF). The Wiener filter is, by construction, the optimal linear estimator that minimizes the mean-squared error between the reconstructed map and the true underlying signal. Our approach leverages a complete statistical description of the sky by incorporating the full matrix of auto- and cross-frequency power spectra for all relevant signal, foreground, and noise components. Instead of attempting to project out the CIB, we explicitly model it as a source of correlated foreground noise. This allows the filter to find the statistically optimal balance between suppressing CIB contamination, preserving the tSZ signal, and minimizing instrumental noise. We apply this MWF to simulated multi-frequency observations representative of the Simons Observatory and Planck, reconstruct the tSZ Compton-$y$ map, and evaluate its performance using a matched-filter cluster detection pipeline. By comparing the purity and completeness of the resulting cluster catalog to one derived from a conventional ILC method, we demonstrate that a full statistical treatment of foregrounds is essential for robustly extracting faint cosmological signals and maximizing the scientific return of next-generation CMB surveys.

\end{document}
                